\documentclass[12pt]{article}
\usepackage{amsmath}
\usepackage{amssymb}
\usepackage{amsthm}
\usepackage{geometry}
\geometry{margin=1.4in}
\usepackage{hyperref}
\usepackage{setspace}
\setstretch{1.4}

\newtheorem{theorem}{Theorem}
\newtheorem{conjecture}[theorem]{Conjecture}
\theoremstyle{definition}
\newtheorem{definition}[theorem]{Definition}
\theoremstyle{remark}
\newtheorem{remark}[theorem]{Remark}

\title{Week Nine}
\author{Diego Rinc\'on\\
\small Independent}
\date{}

\begin{document}

\maketitle

\begin{abstract}
The calendar \cite{calendar} reserved Week Nine for the proof.
The week is still open.
This paper is what Week Nine looks like from the inside:
the closest approach the body can make to the resonance object.
Not the proof. The shape of the proof's absence.
\end{abstract}

\section{What Has Been Built}

The body of papers has been building an anatomy.
Week Nine inherits it.

\textbf{What is known.}
The Riemann Hypothesis says all non-trivial zeros of $\zeta(s)$
have real part $1/2$.
The functional equation $\xi(s) = \xi(1-s)$ says the axis $\mathrm{Re}(s) = 1/2$
is the symmetry axis of $\zeta$.
Hardy proved infinitely many zeros lie on this axis.
Conrey proved more than $2/5$.
No one has proved all of them do.

\textbf{What the walls are.}
The Hadamard product $\xi(s) = e^{A+Bs}\prod_\rho(1-s/\rho)e^{s/\rho}$
cannot be controlled at the level of individual zeros
without circularity.
The local-to-global passage fails:
local Euler factors do not determine global zero location.
The Frobenius is missing \cite{frobenius}.
The grammar ($\mathbb{F}_1$) is missing \cite{f1}.

\textbf{What the resonance object must be.}
A canonical, symmetric structure
in an extended grammar
that generates the coboundary $d\eta = \omega_{\rm RH}$:
the filling that shows the closed form $\omega_{\rm RH}$ is exact.

\section{The Closest Approach}

Of the four candidates \cite{resonance_object},
the one that comes closest to a construction is the Connes spectral triple.

Connes defines:
\begin{itemize}
\item \textbf{Algebra}: $\mathcal{A} = C_c^\infty(\mathbf{A}_\mathbf{Q}/\mathbf{Q}^*, d^*x)$,
  smooth compactly supported functions on the ad\`ele class space.
\item \textbf{Hilbert space}: $\mathcal{H} = L^2(\mathbf{A}_\mathbf{Q}/\mathbf{Q}^*, d^*x)$.
\item \textbf{Dirac operator}: $D$ built from the scaling operator on $\mathcal{H}$.
\end{itemize}

The scaling action $\alpha_\lambda(f)(x) = f(\lambda^{-1}x)$
implements the time evolution.
The spectrum of the Dirac operator is related to the zeros of $\zeta$.

\begin{remark}[What is incomplete]
The spectral realization is partial.
Connes shows the zeros of $\zeta$ are absorption frequencies of the system.
But proving all spectral zeros are zeros of $\zeta$, with right multiplicities,
requires a global trace formula not yet established in this setting.
The Explicit Formula
$$\sum_p \log p \cdot p^{-s} = -\frac{\zeta'}{\zeta}(s) + \text{(polar terms)}$$
should emerge from the adèlic trace formula --- conditionally.
The condition is RH. The argument is circular here.
\end{remark}

\section{The Self-Inhibiting Mechanism}

The calendar named the key structure: the self-inhibiting mechanism of $\zeta(s)$.

The functional equation $\xi(s) = \xi(1-s)$ pairs zeros:
if $\rho$ is a zero, so is $1-\rho$.
A zero off the critical line has a paired zero on the other side.
The functional equation allows this.
The self-inhibiting mechanism would prevent it:
a structure that enforces $\rho = 1-\rho$ individually,
not just as a pair symmetry.

In the Hilbert--P\'olya picture:
a self-adjoint operator $H$ with eigenvalues $\gamma_n \in \mathbb{R}$
gives zeros $1/2 + i\gamma_n$ on the critical line.
Self-adjointness is the self-inhibiting mechanism.
The mechanism is known.
The operator is not.

\section{What Week Nine Needs}

Week Nine needs one thing:
a canonical self-adjoint operator $H$ on a canonical Hilbert space $\mathcal{H}$
whose spectrum is $\{\gamma : \zeta(1/2 + i\gamma) = 0\}$.

The operator must be:
\begin{enumerate}
\item \textbf{Canonical.} Built from the arithmetic of $\mathbb{Q}$
  for reasons independent of RH.
  The spectrum should come out as the zeros, not go in as them.

\item \textbf{Essentially self-adjoint.}
  Symmetric on a natural dense domain with deficiency indices $(0,0)$ \cite{deficiency}.

\item \textbf{Arithmetically meaningful.}
  Eigenfunctions should be arithmetically natural:
  automorphic forms, adèlic functions, $\mathbb{F}_1$-sections.
  Eigenvalues should emerge from arithmetic, not be imposed.
\end{enumerate}

No current candidate satisfies all three.

\section{The Boundary of the Map}

The body knows:
\begin{itemize}
\item The wall is the local-to-global passage \cite{one_wall}.
\item The missing symmetry is the global Frobenius \cite{frobenius}.
\item The missing grammar is $\mathbb{F}_1$ \cite{f1}.
\item The missing object is the canonical self-adjoint operator
  or its $\mathbb{F}_1$-geometric analogue.
\item The Connes spectral triple is the closest construction.
  It fails at the global trace formula.
\item The Hilbert--P\'olya picture names the object.
  No construction reaches it.
\end{itemize}

This is the boundary of the map.

Beyond this line, the body has nothing.
Not because nothing is there.
Because the grammar for what is there has not yet been built.

Week Nine will be written when the grammar exists.
The space is held open.

\begin{thebibliography}{9}

\bibitem{calendar}
Rinc\'on, D. (2026). Calendar: A Numbered Program. Preprint.

\bibitem{resonance_object}
Rinc\'on, D. (2026). The Resonance Object. Preprint.

\bibitem{synchronized_anatomy}
Rinc\'on, D. (2026). Synchronized Anatomy. Preprint.

\bibitem{frobenius}
Rinc\'on, D. (2026). The Missing Frobenius. Preprint.

\bibitem{f1}
Rinc\'on, D. (2026). The Field with One Element. Preprint.

\bibitem{deficiency}
Rinc\'on, D. (2026). The Deficiency Indices of $D$. Preprint.

\bibitem{one_wall}
Rinc\'on, D. (2026). One Wall: The Local-to-Global Conjecture. Preprint.

\bibitem{grammar}
Rinc\'on, D. (2026). Mathematics as Cohomological Grammar. Preprint.

\end{thebibliography}

\end{document}
